\chapter{Complete Examples}

Every example in this chapter is a complete input deck. It includes geometry, properties, boundary conditions, loads where applicable, and the loadcase. The examples can therefore be copied into separate \file{.inp} files and run directly.

\section{Linear static truss}

This first model is a single axial bar in N--mm units. The guide support at the tip removes the two transverse rigid-body modes without constraining axial extension.

\begin{exampledeck}
*MODEL, NAME=LINEAR_TRUSS

*NODE, NSET=NALL
1,    0.0, 0.0, 0.0
2, 1000.0, 0.0, 0.0

*ELEMENT, TYPE=T3, ELSET=BAR
1, 1, 2

*NSET, NAME=FIXED
1
*NSET, NAME=TIP
2

*MATERIAL, NAME=STEEL
*ELASTIC, TYPE=ISOTROPIC
210000.0, 0.3

*TRUSSSECTION, ELSET=BAR, MATERIAL=STEEL
100.0

*SUPPORT, SUPPORT_COLLECTOR=BC
FIXED, 0.0, 0.0, 0.0, NAN, NAN, NAN
TIP,   NAN, 0.0, 0.0, NAN, NAN, NAN

*CLOAD, LOAD_COLLECTOR=AXIAL_LOAD
TIP, 1000.0, 0.0, 0.0, 0.0, 0.0, 0.0

*LOADCASE, TYPE=LINEARSTATIC, NAME=AXIAL_TENSION
*SUPPORTS
BC
*LOADS
AXIAL_LOAD
*SOLVER, DEVICE=CPU, METHOD=DIRECT
*CONSTRAINTMETHOD, TYPE=NULLSPACE
*END
\end{exampledeck}

\section{Topology-aware linear static solid}

The density and orientation fields are defined at root level and referenced by keyword. \cmd{TOPODENSITY} requires \key{FIELD}; the data-line form is not valid. The orientation angles are Euler angles in radians and are optional. Density \(\rho_e=0.8\) with \(p=3\) scales this element's stiffness by \(0.8^3\).

\begin{exampledeck}
*MODEL, NAME=TOPOLOGY_CUBE

*NODE, NSET=NALL
1, 0.0, 0.0, 0.0
2, 0.0, 1.0, 0.0
3, 0.0, 0.0, 1.0
4, 0.0, 1.0, 1.0
5, 1.0, 0.0, 0.0
6, 1.0, 1.0, 0.0
7, 1.0, 0.0, 1.0
8, 1.0, 1.0, 1.0

*ELEMENT, TYPE=C3D8, ELSET=DESIGN_DOMAIN
1, 1, 5, 6, 2, 3, 7, 8, 4

*NSET, NAME=FIXED_FACE
1, 2, 3, 4
*NSET, NAME=LOADED_FACE
5, 6, 7, 8

*MATERIAL, NAME=ORTHOTROPIC_SOLID
*ELASTIC, TYPE=ORTHOTROPIC
100000.0, 50000.0, 50000.0, 20000.0, 20000.0, 25000.0, 0.25, 0.25, 0.25

*SOLIDSECTION, ELSET=DESIGN_DOMAIN, MATERIAL=ORTHOTROPIC_SOLID

*FIELD, NAME=RHO, TYPE=ELEMENT, COLS=1, FILL=ZERO
1, 0.8

*FIELD, NAME=ANGLES, TYPE=ELEMENT, COLS=3, FILL=ZERO
1, 0.0, 0.0, 0.0

*SUPPORT, SUPPORT_COLLECTOR=BC
FIXED_FACE, 0.0, 0.0, 0.0, NAN, NAN, NAN

*CLOAD, LOAD_COLLECTOR=TENSION
LOADED_FACE, 250.0, 0.0, 0.0, 0.0, 0.0, 0.0

*LOADCASE, TYPE=LINEARSTATICTOPO, NAME=TOPO_STATIC
*SUPPORTS
BC
*LOADS
TENSION
*SOLVER, DEVICE=CPU, METHOD=DIRECT
*CONSTRAINTMETHOD, TYPE=NULLSPACE
*TOPODENSITY, FIELD=RHO
*TOPOORIENT, FIELD=ANGLES
*TOPOEXPONENT
3.0
*END
\end{exampledeck}

\section{Shell with explicit ABD and transverse shear stiffness}

The six lines following \cmd{ELASTIC} are the rows of the \(6\times6\) ABD matrix. The final line contains the four entries of \(S\) in row-major order. The example uses a symmetric uncoupled laminate, but the parser accepts all 36 ABD entries independently.

\begin{exampledeck}
*MODEL, NAME=ABD_CANTILEVER

*NODE, NSET=NALL
1,    0.0,   0.0, 0.0
2, 1000.0,   0.0, 0.0
3, 1000.0, 500.0, 0.0
4,    0.0, 500.0, 0.0

*ELEMENT, TYPE=S4, ELSET=PANEL
1, 1, 2, 3, 4

*NSET, NAME=ROOT
1, 4
*NSET, NAME=TIP
2, 3

*MATERIAL, NAME=LAMINATE_ABD
*ELASTIC, TYPE=ABD
1.0e6, 2.5e5, 0.0, 0.0, 0.0, 0.0
2.5e5, 1.0e6, 0.0, 0.0, 0.0, 0.0
0.0, 0.0, 3.0e5, 0.0, 0.0, 0.0
0.0, 0.0, 0.0, 1.0e3, 2.5e2, 0.0
0.0, 0.0, 0.0, 2.5e2, 1.0e3, 0.0
0.0, 0.0, 0.0, 0.0, 0.0, 3.0e2
3.0e5, 0.0, 0.0, 3.0e5

*SHELLSECTION, ELSET=PANEL, MATERIAL=LAMINATE_ABD
1.0

*SUPPORT, SUPPORT_COLLECTOR=BC
ROOT, 0.0, 0.0, 0.0, 0.0, 0.0, 0.0

*CLOAD, LOAD_COLLECTOR=TIP_LOAD
TIP, 0.0, 0.0, -100.0, 0.0, 0.0, 0.0

*LOADCASE, TYPE=LINEARSTATIC, NAME=ABD_BENDING
*SUPPORTS
BC
*LOADS
TIP_LOAD
*SOLVER, DEVICE=CPU, METHOD=DIRECT
*CONSTRAINTMETHOD, TYPE=NULLSPACE
*END
\end{exampledeck}

\section{Nonlinear load control}

The two initially horizontal bars have no linearized vertical stiffness at their midpoint. Zero-row regularization supplies the initial perturbation; after deflection, geometric stiffness carries the load.

\begin{exampledeck}
*MODEL, NAME=TWO_BAR_MECHANISM

*NODE, NSET=NALL
1, -1000.0, 0.0, 0.0
2,     0.0, 0.0, 0.0
3,  1000.0, 0.0, 0.0

*ELEMENT, TYPE=T3, ELSET=BARS
1, 1, 2
2, 2, 3

*NSET, NAME=END_SUPPORTS
1, 3
*NSET, NAME=MIDPOINT
2

*MATERIAL, NAME=STEEL
*ELASTIC, TYPE=ISOTROPIC
210000.0, 0.3

*TRUSSSECTION, ELSET=BARS, MATERIAL=STEEL
100.0

*SUPPORT, SUPPORT_COLLECTOR=BC
END_SUPPORTS, 0.0, 0.0, 0.0, NAN, NAN, NAN
MIDPOINT, NAN, NAN, 0.0, NAN, NAN, NAN

*CLOAD, LOAD_COLLECTOR=MIDPOINT_LOAD
MIDPOINT, 0.0, -1000.0, 0.0, 0.0, 0.0, 0.0

*LOADCASE, TYPE=NONLINEARSTATIC, NAME=LOAD_CONTROL
*SUPPORTS
BC
*LOADS
MIDPOINT_LOAD
*SOLVER, DEVICE=CPU, METHOD=DIRECT
*CONSTRAINTMETHOD, TYPE=NULLSPACE
*NONLINEAR, CONTROL=LOAD, INCREMENTS=20, MAXITER=40, TOL=1e-8, ADAPTIVE=ON, REGULARIZE_ZERO_ROWS=ON, REGULARIZATION_ALPHA=1e-4
*END
\end{exampledeck}

\section{Snap-through with arc-length control}

The 20 kN reference load exceeds the first limit load of the shallow two-bar arch. Arc-length control can follow the equilibrium path while the load factor temporarily decreases.

\begin{exampledeck}
*MODEL, NAME=SNAP_THROUGH_TRUSS

*NODE, NSET=NALL
1, -1000.0,   0.0, 0.0
2,     0.0, 100.0, 0.0
3,  1000.0,   0.0, 0.0

*ELEMENT, TYPE=T3, ELSET=BARS
1, 1, 2
2, 2, 3

*NSET, NAME=END_SUPPORTS
1, 3
*NSET, NAME=APEX
2

*MATERIAL, NAME=STEEL
*ELASTIC, TYPE=ISOTROPIC
210000.0, 0.3

*TRUSSSECTION, ELSET=BARS, MATERIAL=STEEL
100.0

*SUPPORT, SUPPORT_COLLECTOR=BC
END_SUPPORTS, 0.0, 0.0, 0.0, NAN, NAN, NAN
APEX, 0.0, NAN, 0.0, NAN, NAN, NAN

*CLOAD, LOAD_COLLECTOR=APEX_LOAD
APEX, 0.0, -20000.0, 0.0, 0.0, 0.0, 0.0

*LOADCASE, TYPE=NONLINEARSTATIC, NAME=ARC_LENGTH
*SUPPORTS
BC
*LOADS
APEX_LOAD
*SOLVER, DEVICE=CPU, METHOD=DIRECT
*CONSTRAINTMETHOD, TYPE=NULLSPACE
*NONLINEAR, CONTROL=ARC_LENGTH, ARC_LENGTH_PSI=1.0, MAX_INCREMENTS=200, INITIAL_INCREMENT=0.1, MINIMUM_INCREMENT=1e-5, MAXIMUM_INCREMENT=0.1, ADAPTIVE=ON, MAXITER=30, TOL=1e-8, REGULARIZE_ZERO_ROWS=OFF
*END
\end{exampledeck}

\section{Eigenfrequency with a point mass}

The two-element bar has two free axial degrees of freedom; the loadcase requests the lowest mode. Its distributed density and the concentrated tip mass both contribute to the mass matrix.

\begin{exampledeck}
*MODEL, NAME=AXIAL_EIGENFREQUENCY

*NODE, NSET=NALL
1,    0.0, 0.0, 0.0
2,  500.0, 0.0, 0.0
3, 1000.0, 0.0, 0.0

*ELEMENT, TYPE=T3, ELSET=BAR
1, 1, 2
2, 2, 3

*NSET, NAME=FIXED
1
*NSET, NAME=GUIDED
2, 3
*NSET, NAME=TIP
3

*MATERIAL, NAME=STEEL
*ELASTIC, TYPE=ISOTROPIC
210000.0, 0.3
*DENSITY
7.85e-9

*TRUSSSECTION, ELSET=BAR, MATERIAL=STEEL
100.0

*POINTMASS, NSET=TIP
0.01, 0.0, 0.0, 0.0, 0.0, 0.0, 0.0, 0.0, 0.0, 0.0

*SUPPORT, SUPPORT_COLLECTOR=BC
FIXED, 0.0, 0.0, 0.0, NAN, NAN, NAN
GUIDED, NAN, 0.0, 0.0, NAN, NAN, NAN

*LOADCASE, TYPE=EIGENFREQ, NAME=AXIAL_MODE
*SUPPORTS
BC
*NUMEIGENVALUES
1
*END
\end{exampledeck}

\section{Linear buckling of a shallow truss}

The preload places both bars in compression. Linear buckling then evaluates the geometric stiffness of that preloaded state and reports the requested factor and mode.

\begin{exampledeck}
*MODEL, NAME=SHALLOW_TRUSS_BUCKLING

*NODE, NSET=NALL
1, -1000.0,   0.0, 0.0
2,     0.0, 100.0, 0.0
3,  1000.0,   0.0, 0.0

*ELEMENT, TYPE=T3, ELSET=BARS
1, 1, 2
2, 2, 3

*NSET, NAME=ENDS
1, 3
*NSET, NAME=APEX
2

*MATERIAL, NAME=STEEL
*ELASTIC, TYPE=ISOTROPIC
210000.0, 0.3

*TRUSSSECTION, ELSET=BARS, MATERIAL=STEEL
100.0

*SUPPORT, SUPPORT_COLLECTOR=BC
ENDS, 0.0, 0.0, 0.0, NAN, NAN, NAN
APEX, NAN, NAN, 0.0, NAN, NAN, NAN

*CLOAD, LOAD_COLLECTOR=PRELOAD
APEX, 0.0, -1000.0, 0.0, 0.0, 0.0, 0.0

*LOADCASE, TYPE=LINEARBUCKLING, NAME=LIMIT_ESTIMATE
*SUPPORTS
BC
*LOADS
PRELOAD
*SOLVER, DEVICE=CPU, METHOD=DIRECT
*NUMEIGENVALUES
1
*SIGMA
0.0
*END
\end{exampledeck}

\section{Linear transient ramp load}

The amplitude ramps the axial force over 0.1 seconds and then holds it constant. Results are written every ten integration steps.

\begin{exampledeck}
*MODEL, NAME=TRANSIENT_TRUSS

*NODE, NSET=NALL
1,    0.0, 0.0, 0.0
2, 1000.0, 0.0, 0.0

*ELEMENT, TYPE=T3, ELSET=BAR
1, 1, 2

*NSET, NAME=FIXED
1
*NSET, NAME=TIP
2

*MATERIAL, NAME=STEEL
*ELASTIC, TYPE=ISOTROPIC
210000.0, 0.3
*DENSITY
7.85e-9

*TRUSSSECTION, ELSET=BAR, MATERIAL=STEEL
100.0

*SUPPORT, SUPPORT_COLLECTOR=BC
FIXED, 0.0, 0.0, 0.0, NAN, NAN, NAN
TIP,   NAN, 0.0, 0.0, NAN, NAN, NAN

*AMPLITUDE, NAME=RAMP, TYPE=LINEAR
0.0, 0.0
0.1, 1.0
1.0, 1.0

*CLOAD, LOAD_COLLECTOR=DYNAMIC_LOAD, AMPLITUDE=RAMP
TIP, 1000.0, 0.0, 0.0, 0.0, 0.0, 0.0

*LOADCASE, TYPE=LINEARTRANSIENT, NAME=RAMP_RESPONSE
*SUPPORTS
BC
*LOADS
DYNAMIC_LOAD
*SOLVER, DEVICE=CPU, METHOD=DIRECT
*TIME
0.0, 1.0, 0.01
*NEWMARK
0.25, 0.5
*DAMPING, TYPE=RAYLEIGH
0.0, 1e-5
*WRITEEVERY, TYPE=STEPS
10
*END
\end{exampledeck}
